\subsection{Реализация Telegram-бота, HTTP-сервера и HTTP-клиентов}

\subsubsection{Описание}

В качестве первого шага нужно разработать сетевую часть приложения и основную логику бота. 
Задача состоит в написании "<скелета"> бота, добавлении сетевых вызовов согласно OpenAPI-контракту, реализации HTTP-клиентов (для внешних и внутренних вызовов) и планировщика обновлений (проверяющего обновления в каналах).

\textbf{Цель:} Получить работающий \textbf{MVP}. Должна работать отправка сообщений в бота и простая реакция на сообщения пользователей.

\subsubsection{Функциональные требования}

\paragraph{Команды бота}
\begin{itemize}
    \item Бот поддерживает следующие команды (неизвестные команды игнорируются с уведомлением пользователю):

    \begin{description}
        \item[\texttt{/start}] Регистрация пользователя.
        \item[\texttt{/help}] Вывод списка доступных команд.
        \item[\texttt{/track}] Начать отслеживание telegram-канала.
        \item[\texttt{/untrack}] Прекратить отслеживание telegram-канала.
        \item[\texttt{/list}] Показать список отслеживаемых telegram-каналов. Если список пуст — выводить специальное сообщение.
    \end{description}

    \item При попытке добавить telegram-канал, который уже отслеживается, выводится сообщение: \textit{"<Данный telegram-канал уже отслеживается">}.
\end{itemize}

\paragraph*{OpenAPI и Взаимодействие}
\begin{itemize}
    \item Все endpoint'ы соответствуют \href{TODO}{OpenAPI-контракту}.
    \item Общение между приложениями \texttt{bot} и \texttt{scrapper} происходит по HTTP.
\end{itemize}

\paragraph*{Планировщик обновлений} \mbox{} \\
Реализован базовый планировщик:
\begin{itemize}
    \item Бот присылает простейшее уведомление (заглушку) при обнаружении изменений.
    \item Детализация изменений не требуется ("<обновление в канале Y"> достаточно).
    \item Реализована простая версия, не учитывающая множество пользователей.
\end{itemize}

\subsubsection{Нефункциональные требования}

\paragraph*{Безопасность и конфигурация}
\begin{itemize}
    \item Токен авторизации хранится в конфигурационном файле, закрытом для общего доступа.
    \item Используется типобезопасная конфигурация.
\end{itemize}

\paragraph*{HTTP-клиенты и Scrapper} \mbox{} \\
TODO

\paragraph*{Логика взаимодействия (FSM)} \mbox{} \\
Бот реализует концепт машины состояний. Создание подписки идет в формате диалога.

\noindent{Пример диалога (/track):}

\noindent\fbox{
    \parbox{\dimexpr\linewidth-2\fboxsep-2\fboxrule}{
        \ttfamily\small
        > /track @Yandex4Developers \\
        < введите фильтры (опционально) \\
        > development c++ high-load kotlin
    }
}

\paragraph*{Тестирование и Логирование}
\begin{itemize}
    \item В тестах отсутствуют реальные вызовы к API внешних систем (использованы mock-заглушки).
    \item Реализовано структурное логирование (добавление key-value значений к логу вместо зашивания данных в поле message) в коде приложения.
\end{itemize}

\subsubsection{Покрытие тестами}

\begin{itemize}
    \item Корректный парсинг ссылок.
    \item Правильное сохранение данных в репозиторий (ссылка, тэги, фильтры).
    \item Бот кидает ошибку, если команда неизвестна.
    \item \textbf{Happy path}: добавление и удаление ссылок из репозитория.
    \item Тест на добавление дубля канала.
    \item \textbf{Планировщик}: обновление отправляется только пользователям, которые следят за каналом.
    \item Корректная обработка ошибок HTTP-запросов (некорректное тело, код ответа) в клиентах.
    \item Проверка форматирования команды \cmd{/list}.
\end{itemize}
